\section{What Certification Establishes}

If all proof stages are independently accepted, certification establishes:
\begin{itemize}
  \item the enacted recursive schedule was followed;
  \item every parent and child used the declared unit universe;
  \item no permitted connected cut improved the objective at any node; and
  \item the final assignment is the unique canonical BISECT plan.
\end{itemize}

Certification does not establish that the enacted axioms are normatively
correct. It does not optimize partisan symmetry, competitiveness, communities
of interest, racial representation, or every compactness measure. VRA and
constitutional application remain separate legal questions.

Nor does the present implementation show that certified recursion ``wins'' on
nationwide map quality. The precommitted eight-instance suite now runs the same
bounded root instances through exact certified and heuristic selection without
choosing seeds after the result. It is a comparison-mechanics result, not a
representative political or State-scale sample. A defensible superiority study
must run the same instances through certified and heuristic cut selection, report objective
agreement and disagreement, measure runtime and proof size, and evaluate final
plan metrics without selecting favorable states or seeds.

The current conclusion is narrower: relative to heuristic output, solver
metadata, and RPLAN audit alone, certified recursive bisection provides the
strongest model-optimality claim architecture currently implemented in the
project. The bounded comparison shows why proof can matter even when every
heuristic result is feasible: ten of 40 rows sacrifice the model's
first-priority population objective for a smaller boundary. NRS v0.3
demonstrates that the operational recursion, assignment, and
verification layer can run nationally across three Census cycles. Whether the
boundary and canonical proof strength can be delivered at that scale remains
the central open engineering question.
